% Under synthetic corruptions, homogeneous performance decreases consistently with increasing severity. 
% For Gaussian noise, accuracy drops from roughly 83\% at severity level 1 to about 35\% at level 5. 
% The strongest decline occurs between the intermediate and higher severity levels. 
% For pixelation, the degradation is more gradual. 
% Accuracy decreases from approximately 87\% at severity level 1 to around 65\% at severity level 5. 
% Across all severity levels, pixelation leads to a less pronounced performance drop compared to Gaussian noise.

Figures~\ref{fig:appendix_gaussian} and~\ref{fig:appendix_pixelate} show accuracy as corruption severity increases for Gaussian noise and pixelation.

At low corruption levels, both ensembles and the strongest single models perform similarly. 
Differences are small and remain within a narrow range. 

As severity increases, performance begins to diverge. 
Under Gaussian noise, the homogeneous ResNet ensemble follows the degradation pattern of the single ResNet-50 and exhibits a pronounced drop at higher severity levels. 
In contrast, the heterogeneous ensemble shows a flatter decline and maintains higher accuracy under severe corruption. 
The difference becomes most visible at severity levels 4 and 5.

A similar pattern is observed for pixelation. 
At low levels, both ensembles behave comparably. 
With increasing distortion, however, the homogeneous ensemble degrades more strongly, while the heterogeneous ensemble remains more stable. 

Across both corruption types, the main signal lies in the slope of degradation rather than in absolute accuracy at low severity. 
Architectural diversity reduces sensitivity to increasing corruption strength, whereas homogeneous ResNet models tend to fail in a more correlated manner under severe shift.